\section{\textit{Pipeline} en tiempo real con juegos}
\label{sec:pipeline_tiempo_real}

Una vez preparado el modelo de clasificación, el siguiente paso del proyecto consistió en integrar dicho modelo dentro de un sistema interactivo en tiempo real. El objetivo principal de esta etapa fue comprobar la viabilidad práctica del uso de señales EEG para el control de aplicaciones dinámicas, concretamente videojuegos sencillos diseñados específicamente para el experimento.

A diferencia de un entorno \textit{offline}, donde las predicciones se pueden dividir y procesar en \textit{trials} completos y procesados, en un entorno en tiempo real se plantean nuevos problemas relacionados con la latencia, la estabilidad de las predicciones y la continuidad temporal de las decisiones tomadas por el modelo. Debido a ello, fue necesario diseñar un \textit{pipeline} capaz de recibir muestras EEG continuamente, procesarlas mediante el modelo entrenado y transformar las predicciones en acciones útiles dentro del juego.

Se decidió usar un sistema modular con una arquitectura definida específicamente para dividir las responsabilidades, que puedan trabajar en paralelo y permitir cambiar bloques si fuera preciso.

Con objetivo de analizar las salidas del modelo en un entorno de tiempo real, se planteó un sistema ``juguete'' que recogía las predicciones de una grabación \textit{offline} y las enviaba al sistema, emulando ser lecturas en tiempo real. Analizando los resultados de este experimento, se llegó a la siguiente conclusión: era necesario aplicar, tanto un ``filtro de decisión'' como un ``filtro de ayuda'' (similar a un \textit{aimbot}) para poner orden en el caos que suponían las predicciones crudas del modelo.

Además, se diseñaron y desarrollaron varios juegos experimentales que sirvieron para comprobar la capacidad de integración de nuestro sistema en un entorno de tiempo real. Estos juegos incorporan mecanismos de asistencia adicional, orientados a reducir los errores derivados del propio ruido de las señales EEG, de las fluctuaciones del modelo de clasificación y de la experiencia del usuario con \textit{motor imagery}.

\subsection{Arquitectura del \textit{pipeline}}
\label{subsec:arquitectura}

\subsection{Filtros de decisión}
\label{subsec:filtros_decision}

El objetivo principal de estos ``filtros de decisión'' es hacer más consistente el flujo de predicciones que se recibe del modelo. Esto se puede conseguir de distintas formas, entre ellas: reducir la cantidad de predicciones agrupándolas, usar mecanismos para ignorar \textit{outliers} (predicciones que se salen de lo ``normal'' en ese periodo temporal), suavizando los cambios de una predicción a otra (eliminando saltos).

Se plantearon y testearon varias opciones:

\begin{itemize}
    \item \textbf{Filtro por mayoría}: Consiste en, dentro de una ventana de tiempo determinado (el último segundo, los últimos 4 segundos\dots), contar el número de veces que se ha recibido la predicción ``derecha'' y cuántas veces se ha predicho ``izquierda''. La decisión final se obtiene seleccionando la clase con mayor frecuencia de aparición en la ventana seleccionada. Este método proporciona una estrategia sencilla frente a predicciones aisladas erróneas.
    \item \textbf{\textit{Score} ponderado}:  Para cada predicción generada por el modelo, se calcula un valor de confianza a partir de la diferencia absoluta entre las probabilidades asociadas a ambas clases: $\textit{Score} = abs(P_derecha - P_izquierda)$. Cuanto mayor sea dicha diferencia, mayor será la confianza atribuida a la predicción. Posteriormente, cada predicción se pondera según este \textit{score}, asignando valores positivos a una clase (\textit{right\_hand}) y negativos a la otra (\textit{left\_hand}). Todo esto se hace con el objetivo de obtener una media ponderada que tenga en cuenta tanto la dirección predicha como la seguridad del modelo.
    \item \textbf{Media de probabilidades}: Con este método se promedian independientemente las probabilidades asociadas a cada clase obtenidas en todas las predicciones de una ventana temporal determinada. La decisión final corresponde a la clase cuya probabilidad media resulte mayor. Este enfoque permite aprovechar directamente la información probabilística proporcionada por el modelo y trata de mejorar el método de la mayoría.
    \item \textbf{Filtro con fuga}: Este se basa en un mecanismo acumulativo inspirado en sistemas dinámicos con disminución temporal. A partir de cada predicción se incrementa un acumulador asociado a la clase correspondiente. De forma periódica, el valor acumulado disminuye mediante un factor de fuga. Cuando el acumulador de una clase supera un determinado umbral, se genera la decisión correspondiente.
    \item \textbf{\textit{Exponential smoothing}}: Este método aplica un filtrado temporal recursivo sobre las probabilidades generadas por el modelo, combinando la predicción actual con el histórico de predicciones anteriores. Para ello, se calcula una media exponencialmente ponderada en la que las observaciones más recientes tienen un mayor peso que las más antiguas. El grado de influencia de las nuevas predicciones se controla mediante un parámetro de suavizado. De esta forma, el sistema reduce fluctuaciones bruscas y variaciones instantáneas provocadas por el ruido de las señales EEG, manteniendo al mismo tiempo capacidad de adaptación frente a cambios sostenidos en la intención del usuario. La decisión final se obtiene seleccionando la clase cuya probabilidad suavizada resulte mayor.
\end{itemize}

\subsection{Juegos propuestos para el experimento}
\label{subsec:juegos}

Teniendo en cuenta los resultados obtenidos del modelo y todas las pruebas anteriores, se decidió trabajar con tan solo dos predicciones: mano derecha (\textit{right\_hand}) y mano izquierda (\textit{left\_hand}). Por ello, los juegos propuestos para el experimento debían tener como máximo dos controles.

Con esto en mente, se plantearon dos juegos que permitían poner a prueba la precisión y el tiempo de respuesta del sistema, además de la experiencia del usuario: \textit{wave runner} y \textit{project arrows}.

Para no perder tiempo extra en la implementación de los juegos, se decidió diseñarlos y pedir a un LLM, por medio de un prompt detallado, que proporcionara la implementación. Una vez generada, se revisó el código, se comprobó que cumplían todos los requerimientos y se hicieron algunos ajustes manuales.

\subsubsection{\textit{Endless Wave Runner}}

Para el diseño de este juego se tomó como referencia la nave \textit{``wave''} del juego \textit{Geometry Dash}, de ahí el nombre. En la versión del \textit{Geometry Dash}, la nave avanza automáticamente de izquierda a derecha (realmente la nave permanece estática y es el escenario el que se mueve para dar la sensación de avance) y el jugador controla si se mueve en diagonal hacia arriba o hacia abajo, cambiando entre las dos opciones con una pulsación en la pantalla. El objetivo del juego es avanzar sin que la nave se choque, hasta terminar la fase. 

A diferencia del original, en esta versión el jugador avanza automáticamente de abajo hacia arriba, puede controlar si moverse en diagonal hacia la izquierda o hacia la derecha y la pantalla es ``infinita'' (haciendo que el objetivo sea aguantar lo máximo posible). Se consideró apropiado este juego debido a su simplicidad de controles y jugabilidad: el usuario tiene una tarea clara y sencilla que puede controlar con lo que ofrece nuestro sistema de forma intuitiva. 

A la hora de diseñar el juego se tuvo en cuenta que este fuera configurable, para adaptar las velocidades y la dificultad a las limitaciones del casco EEG y de nuestro propio sistema. Se puede adaptar fácilmente, con tan solo cambiar los atributos de la clase: la velocidad en ambos ejes, la generación de obstáculos, el ancho inicial entre obstáculos, la distancia mínima entre obstáculos, el aumento de velocidad con el tiempo, la disminución del ancho del ``túnel'' (ancho entre obstáculos) y el \textit{respawn} automático, entre otros.

Como se ha mencionado anteriormente, se hizo uso de un LLM para la implementación, concretamente Gemini. El prompt utilizado para la generación fue el siguiente:

\begin{lstlisting}[
    breaklines=true,
    breakatwhitespace=true,
    basicstyle=\ttfamily\small,
    frame=single,
    columns=fullflexible,
    keepspaces=true,
    breakindent=0pt
]
Quiero que me hagas un juego basándote en la flecha del geometry dash. En el geometry dash existe un vehículo que es una flecha, este solo puede ir en diagonal hacia arriba derecha o hacia abajo derecha y cada vez q se pulsa en la pantalla cambia entre las dos opciones. Vamos a hacer una modificación a este vehículo para convertirlo en un juego nuevo y aislado:

En vez de avanzar de izquierda a derecha, nuestra flecha avanzará de abajo a arriba. De forma que la flecha tendrá dos opciones: ir hacia arriba izquierda en diagonal o hacia arriba derecha en diagonal. Y el cambio de los controles se hará mediante las flechas: al pulsar la flecha izquierda, nuestra flecha nave apuntará y se moverá hacia arriba izquierda; y al pulsar la flecha derecha nuestra nave apuntará y se moverá hacia arriba derecha. Debes añadir más funcionalidades para hacer pruebas: si se pulsa la flecha de arriba, nuestra flecha nave deberá apuntar directamente hacia arriba recto. Debe de haber una constante en el script q me permita activar/desactivar los obstáculos, si se generan obstáculos habrá a la izquierda y a la derecha de la pantalla pinchos sobresaliendo de distintos tamaños, algunos muy grandes, otros más pequeños (igual q en las zonas del geometry dash q se juegan con esta nave), si el jugador toca un obstáculo, pierde; en caso de q los obstáculos estén desactivados, la única forma de perder será saliéndose de la pantalla por los laterales (izquierdo y derecho). 
Deben existir constantes en el script para modificar la velocidad del juego, el avance por cada tick de la nave, el tiempo q pasa entre ``tick'' y ``tick'', etc. El avance de la nave debe de ser similar a un scroll, la nave realmente se quedará estática en su coordenada y, solo modificará su coordenada en x y su rotación, y será el propio mapa (el fondo (con algún dibujo que no maree para ver q avanza el juego) y, en caso de q los haya, los obstáculos) el que baje a la velocidad determinada por las variables. También deberá haber unas constantes para controlar la ``dificultad'', de forma q si se generan obstáculos, estos podrán dejar más o menos espacio para que el jugador maniobre en base a esa dificultad. Tema puntos: elige la forma q quieras de contarlos pero q haya puntos y q cuanto más tiempo sobrevivas, más puntos tengas. Debe también existir una constante que determine si hay ``respawn automático'' o no, si lo hay, la flecha al chocarse deberá reposicionarse en el ``lugar más seguro'' dentro de su misma coordenada y, dejando la mayor cantidad de hueco tanto a la izquierda como a la derecha sin q haya obstaculo; sin embargo si no hay respawn automático, el juego terminará, pondrá una pantalla ``game over'' y le indicará los puntos al jugador, además de una opción de volver a jugar con la tecla ``esc''. Por último, deberán existir otras constantes q determinen si el juego va a ir aumentando su dificultad o no: una constante q determine si se va a ir incrementando la velocidad, otra que diga cuanto se aumenta cada vez, y otra que determine cada cuanto se aumenta; otra que determine si se va a ir disminuyendo el espacio que hay para pasar entre los obstáculos, otra que determine cuánto disminuye cada vez y otra que determine cada cuánto se disminuye. 

Notas adicionales: Los obstáculos y el fondo deben de ir a la misma velocidad.

Todas las constantes definidas deberán tener un nombre significativo para que las pueda entender sin ninguna explicación a parte. También deberás explicar con un breve comentario al principio de cada función q hace la misma y, si es muy compleja, detallar de forma más detallada q algoritmo sigue para q podamos interpretarlo.

IMPORTANTE: siempre debe de haber un hueco lo suficientemente grande como para que la nave pase, el juego debe de ser posible de jugar, no puede haber en ningún momento una barrera que sea todo obstáculo ni dejar tan poco espacio que el jugador no pueda maniobrar.

En python con pygame. El juego debe estar contenido en una clase llamada EndlessWaveRunner, que se pueda utilizar desde fuera para llamar al juego.
\end{lstlisting}

\subsubsection{\textit{Project Arrows}}

La inspiración para este juego fue \textit{``Project Rhombus''}, de \textit{Steam}. En el orignial, el jugador controla un rombo situado en mitad de la pantalla, este rombo tiene una flecha en su interior que apunta hacia arriba, abajo, izquierda o derecha de forma excluyente. Con las cuatro flechas del teclado, se controla hacia dónde apunta la flecha del rombo. A medida que avanza el juego, aparecen flechas por una de las cuatro direcciones posibles en el borde de la pantalla y avanzan hacia el centro. El objetivo de el jugador es defenderse, apuntando la flecha del rombo hacia la flecha enemiga que le vaya a alcanzar, si una flecha enemiga llega al rombo y este no está apuntando hacia ella, el jugador pierde.

En \textit{``Project Arrows''} se ha adaptado a nuestro sistema quitando los controles de arriba y abajo, dejando únicamente izquierda y derecha. Se planteó este juego por los mismos motivos que el anterior.

Al igual que el anterior juego, se contempló que este fuera configurable y adaptable a nuestro sistema. Pudiendo modificar atributos como: la velocidad de las flechas enemigas, el tiempo de aparición de las mismas, el aumento de velocidad de estas, la posibilidad de reiniciar el juego automáticamente al morir\dots

Se contempló la posibilidad de añadir un tipo de flecha especial que cambia de lado a mitad de camino, se llegó a implementar pero finalmente se descartó por resultar demasiado complicado.

El prompt utilizado para la generación mediante el LLM Gemini fue:

\begin{lstlisting}[
    breaklines=true,
    breakatwhitespace=true,
    basicstyle=\ttfamily\small,
    frame=single,
    columns=fullflexible,
    keepspaces=true,
    breakindent=0pt
]
Quiero un juego que se base en el project rhombus, en caso de que no sepas cómo funciona te hago un resumen: en el centro está el ``personaje'' jugable, es una flecha que apunta arriba, abajo, a la izquierda O a la derecha (solo a una de esas 4 a la vez). Tiene música de fondo para acompañar los ritmos y van llegando flechas de esos 4 lados (desde fuera de la pantalla hacia adentro), si te toca una flecha y estás ``mirando'' (apuntando con tu personaje flecha) a otro lado, esta te mata y pierdes. El objetivo del juego es aguantar lo máximo posible y la puntuación se mide en segundos de la forma: 00.00 (siendo la parte decimal fracciones de un segundo). Ahora, quiero que realicemos una modificación para hacer un juego jugable con tan solo dos flechas (la izquierda y la derecha), lo q debes hacer es replicar el juego, hacer q solo podamos ``apuntar'' hacia la derecha o hacia la izquierda y que las flechas ``a bloquear'' también vengan únicamente de esos dos lados (desde fuera de la pantalla hacia adentro). Vamos a entrar en varios detalles importantes q debes añadir/controlar: debe haber constantes para controlar distintos aspectos del juego: una constante para controlar la velocidad de las flechas que vienen, otra para indicar si se va a ir aumentando esa velocidad con el paso del tiempo, otra que indique cuánto aumenta la velocidad por cada 0.01s en caso de que la de aumento sea True, otra constante para controlar si se generan flechas ``invertidas'' o no (en caso negativo solo se generarán flechas normales, que apuntan hacia el centro desde el principio).

Nota: de normal las flechas que vienen, apuntan hacia el centro (te apuntan ``a ti'' como jugador), sin embargo pueden existir flechas ``invertidas'', estas son flechas que vienen dadas la vuelta (por ejemplo, si vienen de el lado derecho, avanzan hacia el izquierdo pq quieren ir al centro pero apuntan hacia la derecha, van ``al revés'') estas, cuando están relativamente cerca del jugador, cambian de lado (si estaba en la izquierda, se mueve a la derecha y vice versa), haciendo q ahora la flecha apunte hacia el centro, y avanzan igualmente hacia el centro. son flechas diseñadas para confundir, son de otro color distintivo para verlas más fácil y no son muy complicadas de tratar dependiendo de la velocidad de las flechas. Importante para la jugabilidad y que el jugador no se pierda: las flechas invertidas deben cambiar de lado a una distancia prudencial y de una forma VISUAL (no quiero que se teletransporten, quiero que se desplacen visualmente, dejando tiempo de reacción, de un lado a otro) para que le jugador no las pierda de vista, y deben de mantener su color distintivo aunq ahora estén donde deban estar. 

si el jugador pulsa esc, el juego deberá pausarse y marcar en una pantalla que el juego está pausado, cuanto tiempo lleva sobrevivido hasta ahora el jugador y ``esc para continuar''. Si el jugador muere, dependerá de una constante si se reinicia la puntuación a 0 y se sigue con el juego (respawn automático = True) o si se para la ejecución, se cambia a una pantalla de game over, con el tiempo sobrevivido y un ``esc para reiniciar''

Notas adicionales: importante q la estética sea agradable. 

En python con pygame. El juego debe estar contenido en una clase llamada ProjectArrows, que se pueda utilizar desde fuera para llamar al juego.
\end{lstlisting}

\subsection{Filtros de ayuda integrados en los juegos}
\label{subsec:filtros_ayuda}

Como ha sido mencionado anteriormente, se diseñaron mecanismos de ayuda para cada juego que permiten reducir los fallos causados por las señales EEG, el sistema o el propio usuario haciendo \textit{motor imagery}. Estos mecanismos, al depender directamente del estado del juego y de las normas del mismo, fueron implementados directamente en el \textit{script} de cada juego. 

El planteamiento seguido fue el siguiente:

\begin{enumerate}
    \item Se diseñó una solución óptima capaz de decidir para cada instante una solución que evite perder.
    \item Se adaptó la solución para que, en vez de tomar la decisión directamente, tuviera en cuenta la salida del modelo: en base a una constante definida (``SEVERIDAD\_AIMBOT'') la ayuda decide si actuar siguiendo la predicción, o corregir la acción.
    \item Se modificó la ayuda ya que, en los casos que se usaba con su severidad al máximo (``SEVERIDAD\_AIMBOT'' = 1, no teniendo en cuenta las predicciones del modelo), resultaba excesivamente antinatural y notorio. La modificación implicó dejar márgen desde que el \textit{aimbot} ``ve'' qué acción debería realizar hasta que realmente la hace. 
\end{enumerate}

Las estrategias resultantes fueron:

\begin{itemize}
    \item \textbf{Ayuda \textit{Wave Runner}}: La implementación de la asistencia en este juego se basa en un cálculo geométrico que evalúa el nivel de peligro de la nave para cada instante. El algoritmo exacto se divide en:
    \begin{enumerate}
        \item \textbf{Cálculo de los límites laterales:} En cada \textit{frame}, el sistema identifica el segmento del túnel en el que se encuentra la nave exactamente (basándose en su coordenada Y, que es fija). Calcula los bordes izquierdo (\textit{left\_bound}) y derecho (\textit{right\_bound}) disponibles en esa altura.
        \item \textbf{División del espacio:} A partir de los bordes calculados, se obtiene el centro ideal del túnel para esa Y concreta. El espacio disponible se divide manteniendo un margen alrededor del centro (\textit{tunnel\_center\_left} y \textit{tunnel\_center\_right}). Si la nave navega dentro de esta porción central, el sistema lo considera zona segura y mantiene el rumbo. Esta división se realizó con el objetivo conseguir que el \textit{aimbot} resulte natural.
        \item \textbf{Cálculo dinámico del umbral:} Cuando la nave se sale de la porción central y se acerca a una de las paredes, se calcula el umbral. El algoritmo lo calcula en base a la constante de severidad (\texttt{SEVERIDAD\_AIMBOT}) y en base a la distancia al borde. Se calcula la distancia de la nave al borde más cercano y se pondera de forma que cuanto más cerca esté del peligro, más alto será el umbral. El umbral resultante será como mínimo \texttt{SEVERIDAD\_AIMBOT - 1}.
        \item \textbf{Decisión e intervención:} Como ya se ha mencionado, la ayuda no actua de forma estricta, compara el umbral dinámico con las probabilidades suavizadas que emite el modelo de IA (como \texttt{p\_right\_smooth}). Por ejemplo, si la nave está a punto de chocar por la izquierda y la probabilidad suavizada de ir a la derecha supera el umbral dinámico exigido, la ayuda toma la acción de ir a la derecha. En caso contrario, se respeta la predicción cruda del usuario (\texttt{string\_to\_action}).
    \end{enumerate}
    De esta forma, se logra el objetivo del diseño: el \textit{aimbot} no actúa de forma intrusiva, sino que suaviza la jugabilidad dejando margen de reacción que varía en función de cuán crítica sea la situación y de la severidad especificada.
    \item \textbf{Ayuda \textit{Project Arrows}}: La ayuda para este segundo juego se centra en la evaluación de la amenaza más cercana. El algoritmo procesa el estado del juego siguiendo estos pasos:
    \begin{itemize}
        \item \textbf{Identificación de la amenaza principal:} En cada llamada, el sistema evalúa todas las flechas enemigas activas en pantalla y calcula su distancia respecto al centro del jugador (\textit{CENTER\_X}). Esto permite aislar y seleccionar únicamente la flecha enemiga más cercana, priorizando el peligro más inmediato.
        \item \textbf{División del espacio:} Igual que en el anterior caso, se toman medidas para evitar que la ayuda actúe de forma prematura y vuelva el juego excesivamente automatizado. El sistema solo contempla la posibilidad de intervenir si la amenaza más cercana ha sobrepasado un límite específico.
        \item \textbf{Umbral e intervención:} A diferencia del cálculo dinámico tratado anteriormente, aquí se establece un umbral de intervención directo basado en la constante de configuración (\texttt{SEVERIDAD\_AIMBOT}). Si la amenaza cruza la zona crítica, la asistencia comprueba las probabilidades del modelo. Si, por ejemplo, el impacto inminente proviene de la izquierda y la probabilidad suavizada de ir hacia la izquierda (\texttt{p\_left\_smooth}) es mayor o igual que el umbral, el sistema asume la intención del usuario e impone la acción de bloqueo. Si el umbral no se alcanza, la asistencia no interviene y se ejecuta la predicción cruda del modelo (\texttt{string\_to\_action}).
    \end{itemize}
    Con esta estrategia, se permite que el jugador tenga tiempo de ver la amenaza, identificarla y reaccionar en consecuencia.
\end{itemize}

\subsection{\textit{Benchmark} visual}
\label{subsec:benchmark_visual}

Además de los juegos definidos, se diseñó un módulo de feedback visual diseñado para tener una referencia de la destreza del usuario y de la latencia del sistema. 

Dicho componente, al ejecutarse muestra una ventana en negro con dos flechas, una apuntando hacia la izquierda y la otra hacia la derecha. Está diseñado de forma que se puede cambiar el comportamiento con tan solo llamar a una función o a otra. Sin embargo, las funciones que nos interesan son las que permiten la funcionalidad que se mostrará en las pruebas que realizaremos más adelante. 

Cada flecha actua como un contenedor, si llega una predicción se comprueba que el contenedor opuesto esté vacío: en caso de que lo esté, se suma un valor definido (\texttt{self.step}) al contenedor correspondiente a la predicción; en caso contrario, se le resta el doble de la cantidad que se hubiera sumado (\texttt{self.step * 2}). Por otro lado, en caso de que reciba \textit{``rest''}, se llamará a otra función que simplemente resta la misma cantidad (\texttt{self.step}) al contenedor que no esté vacío (en caso de que alguno de los dos tenga contenido). De esta forma, solo podrá rellenarse parcial o totalmente un contenedor a la vez, mientras el otro se mantiene vacío.

La implementación, de la misma forma que con los juegos, se realizó utilizando un LLM (Gemini). Se le proporcionó el siguiente prompt:

\begin{lstlisting}[
    breaklines=true,
    breakatwhitespace=true,
    basicstyle=\ttfamily\small,
    frame=single,
    columns=fullflexible,
    keepspaces=true,
    breakindent=0pt
]
Quiero hacer un script en python que actúe de componente, se pueda importar desde otros scripts e interactuar con este. La funcionalidad del script debe ser la siguiente: cuando se ejecute debe de abrirse una ventana con pygame que muestre con fondo negro, dos flechas de color azul oscuro (agradable a la vista) vacías, es decir solo debe de mostrar el contorno. las flechas deben de estar dispuestas una al lado de la otra, centradas de altura, la de la izquierda debe de apuntar hacia la izquierda, dejar un hueco no muy grande pero estético y debe de disponerse la otra flecha apuntando a la derecha. Ahora bien, eso nada más iniciarlo, pero el "relleno" de las flechas debe de ser en base a un porcentaje que el script debe de tener dentro del objeto como atributo de inicialización, la flecha izquierda tendrá un porcentaje y la derecha otro, así que dos variables. si el porcentaje de la izquierda es 0, la flecha izquierda debe de estar vacía excepto el contorno, si el porcentaje es un 50% debe de estar llena hasta la mitad, si es 100% debe de estar llena por completo, si es un porcentaje más raro y preciso como 64% la flecha deberá de estar rellena un 64%; lo mismo para la flecha derecha pero rellenándose en el sentido contrario. Las flechas se rellenarán desde el centro hacia fuera, es decir si ambas tienen un 1%, será la parte más cercana al centro la que se verá un poco rellena, si están al 70% será la parte más cercana al centro (el 70% de la flecha) la que se rellenará, y el otro 30% más lejano se quedará vacío excepto por el borde. El relleno debe de ser color azul también pero algo más claro y brillante (igualmente que sea agradable a la vista). y el objeto dentro del script debe de tener funciones para poder cambiar el valor de ambos porcentajes de forma independiente o a la vez, al igual que debe de tener una forma de resetear el objeto (poniendo ambos porcentajes a 0) y de consultar los porcentajes de las flechas. El color quiero poder cambiarlo yo manualmente en el código así que ponme claro dónde se define, y el tamaño de ventana debe de ser no muy grande, y modificable, de forma que si se ajusta el tamaño de la ventana las flechas ajusten su tamaño en proporción de forma acorde y sigan centradas de altura y bien dispuestas (dejando un hueco entre ellas en el medio) de forma correcta.
\end{lstlisting}

Partiendo de la base que generó Gemini, fue necesario hacer una pequeña modificación para hacer las flechas completamente simétricas y se implementaron a mano los métodos descritos en el párrafo anterior al prompt.

\subsection{Integración en la aplicación}
\label{subsec:integracion_app}